\chapter{먼저 약속하고 시험 보기}

\section{텍사스의 명사수}

옛날 텍사스에 총을 아주 잘 쏜다고 소문난 사람이 있었대요.
그 사람 헛간 벽에는 과녁이 여러 개 그려져 있었고, 총알 자국이 모두 과녁 한가운데에 박혀 있었어요.
비결이 뭐였을까요?
그 사람은 먼저 벽에 총을 아무렇게나 쏘고, \emph{나중에} 총알 자국 둘레에 과녁을 그렸어요.

웃긴 이야기지만, 과학자들도 이런 실수를 해요.
결과를 많이 들여다보면 그중 몇 개는 우연히 멋져 보이기 마련이에요.
그 멋진 것 둘레에 과녁을 그리고 “맞혔다!”고 말하면 안 되겠지요.
이 실수를 \textbf{텍사스 명사수의 오류}라고 불러요.

13장 끝에서 제가 찾은 것, “동전 랭크가 뒷면 랭크보다 새 허락을 잘 맞힌다”는 바로 총알 자국 둘레에 그린 과녁일 수 있었어요.
결과를 보고 나서 찾았으니까요.
이것이 진짜인지 알려면 방법은 하나뿐이에요.
\emph{과녁을 먼저 그리고, 아무도 본 적 없는 새 벽에 다시 쏘는 것}.

\section{과녁을 먼저 그렸어요}

저는 2026년 9월 28일에 계획서를 썼어요.
무엇을 맞힐지, 누구를 줄 세울지, 몇 점이면 통과인지를 모두 적었어요.
그리고 10월 1일, 이 계획서와 설정 파일을 공개 저장소에 올려 상태를 “등록됨”으로 바꾸었어요.
이렇게 결과를 보기 전에 계획을 공개하고 못 바꾸게 하는 것을 \textbf{사전 등록}이라고 해요.
날짜 도장이 찍힌 타임캡슐을 묻는 것과 비슷해요.

새 벽은 \textbf{2026년 6월부터 8월까지}의 기록이었어요.
계획서를 쓸 때 이 기록은 아직 아무도 꺼내 보지 않았어요.
저는 기록을 꺼내 오는 프로그램이 등록 전에는 아예 움직이지 않도록 만들어 두었어요.
등록을 마친 뒤, 같은 날 6\textasciitilde{}8월 기록을 꺼내 왔어요.
허락 줄 163,553개, 송금 줄 1,013,023개, GMX 거래 닫기 144,152개였어요.

이번에는 \textbf{5월 31일}이 동결일이에요.
10장 달력에서 “14장”이라고 적힌 줄이에요.
5월 31일까지의 모든 기록으로 점수를 매겨 봉투에 넣고, 6\textasciitilde{}8월에 일어난 일로 채점해요.

\section{세 가지 약속}

계획서에는 동전 랭크와 뒷면 랭크를 비교하는 약속이 세 개 적혀 있었어요.
뒷면 랭크는 9장에서 본, 송금 지도만 걷는 동전 랭크예요.
\emph{세 개를 모두} 지켜야 “동전 랭크가 송금 지도만 걷는 것보다 낫다”고 말할 수 있다고 미리 정했어요.

\begin{center}
\begin{tcolorbox}[colback=white, colframe=inkblue, boxrule=0.8pt, arc=2mm, width=0.95\linewidth]
\small
\textbf{약속 1} \quad 스펜더 1,964개의 새 허락 맞히기에서 동전 랭크가 뒷면 랭크보다 \emph{낫다}.\par\smallskip
\textbf{약속 2} \quad 같은 스펜더들의 새 송금자 맞히기에서 동전 랭크가 뒷면 랭크보다 \emph{0.02 넘게 못하지는 않다}.\par\smallskip
\textbf{약속 3} \quad GMX 거래자들의 “청산 없이 닫은 비율” 맞히기에서 동전 랭크가 뒷면 랭크보다 \emph{0.02 넘게 못하지는 않다}.
\end{tcolorbox}
\end{center}

약속 3은 새로 넣은 시험이에요.
6\textasciitilde{}8월에 거래를 세 번 이상 닫은 거래자마다, 닫은 거래 가운데 청산당하지 않은 비율을 계산해요.
허락 지도에도 송금 지도에도 없는, GMX에서만 나오는 기록이에요.
돈을 잃어 강제로 끝나는 일이 적은 사람일수록 믿을 만하다고 볼 수 있으니까요.
거래자 5,151명 가운데 1,402명이 이 시험에 들어왔어요.

\section{“못하지는 않다”는 어떻게 판정할까요?}

약속 2와 약속 3은 “이긴다”가 아니라 “크게 지지는 않는다”는 약속이에요.
판정에는 11장의 다시 뽑기로 구한 95\% 범위를 써요.

\begin{itemize}
\item “낫다”(약속 1): 차이의 95\% 범위가 모두 0보다 위에 있어야 해요. 운이 나빠도 이긴다는 뜻이에요.
\item “0.02 넘게 못하지는 않다”(약속 2, 약속 3): 차이의 95\% 범위의 아래 끝이 $-0.02$보다 위에 있어야 해요.
      운이 나쁜 쪽으로 봐도 0.02보다 많이 지지는 않는다는 뜻이에요.
\end{itemize}

\section{채점 결과}

\begin{figure}[h]
\centering
\begin{tikzpicture}[x=26cm, y=0.9cm]
  % 축
  \draw[-{Stealth}, line width=0.7pt] (-0.05,0) -- (0.245,0);
  \foreach \v/\l in {-0.05/$-0.05$, 0/$0$, 0.05/$0.05$, 0.1/$0.10$, 0.15/$0.15$, 0.2/$0.20$}
    \draw (\v,0.08) -- (\v,-0.08) node[below, font=\scriptsize] {\l};
  % 기준선
  \draw[dashed, color=roseline, line width=0.9pt] (-0.02,3.4) -- (-0.02,0);
  \node[font=\scriptsize, color=roseline, anchor=south] at (-0.02,3.4) {$-0.02$};
  \draw[dotted, color=greycol, line width=0.8pt] (0,3.4) -- (0,0.1);
  % 약속 1
  \draw[line width=2.2pt, color=leafline] (0.132,2.7) -- (0.225,2.7);
  \fill[color=leafline] (0.179,2.7) circle[radius=2.6pt];
  \node[font=\scriptsize, anchor=south] at (0.179,2.8) {$+0.179$ \quad [0.132, 0.225]};
  \node[font=\scriptsize, anchor=west] at (0.03,2.7) {\textbf{약속 1} \textcolor{leafline}{\textbf{통과}}};
  % 약속 2
  \draw[line width=2.2pt, color=roseline] (-0.025,1.8) -- (0.003,1.8);
  \fill[color=roseline] (-0.011,1.8) circle[radius=2.6pt];
  \node[font=\scriptsize, anchor=west] at (0.015,1.8) {\textbf{약속 2} \textcolor{roseline}{\textbf{실패}} \quad $-0.011$ \quad [$-0.025$, 0.003]};
  % 약속 3
  \draw[line width=2.2pt, color=leafline] (0.011,0.9) -- (0.024,0.9);
  \fill[color=leafline] (0.018,0.9) circle[radius=2.6pt];
  \node[font=\scriptsize, anchor=west] at (0.035,0.9) {\textbf{약속 3} \textcolor{leafline}{\textbf{통과}} \quad $+0.018$ \quad [0.011, 0.024]};
\end{tikzpicture}
\caption{동전 랭크에서 뒷면 랭크를 뺀 $\tau$의 차이와 95\% 범위. 빨간 점선은 “0.02까지는 져도 된다”는 선이에요.}
\end{figure}

\begin{itemize}
\item \textbf{약속 1 통과.} 새 허락 맞히기에서 동전 랭크가 뒷면 랭크보다 0.179 높았고, 95\% 범위는 0.132에서 0.225였어요.
      봄에 눈에 띈 것이 새 벽에서도 다시 나타났어요.
\item \textbf{약속 2 실패.} 새 송금자 맞히기에서 동전 랭크가 뒷면 랭크보다 0.011 낮았어요.
      그 자체는 작은 차이지만, 95\% 범위의 아래 끝이 $-0.025$였어요.
      $-0.02$보다 0.005만큼 더 내려갔어요.
      아주 아슬아슬하게, 약속을 지키지 못했어요.
\item \textbf{약속 3 통과.} 청산 없이 닫은 비율 맞히기에서 동전 랭크가 뒷면 랭크보다 0.018 높았고, 범위는 0.011에서 0.024였어요.
\end{itemize}

세 약속을 모두 지켜야 했는데 하나를 지키지 못했어요.
그래서 판정은 \textbf{미충족}이에요.
저는 논문에서 “동전 랭크가 송금 지도만 걷는 것보다 낫다”고 주장하지 않아요.

\section{아깝지 않냐고요?}

솔직히 말하면 아까웠어요.
0.005 차이였으니까요.
“범위를 조금만 다르게 잡았다면”, “0.02 대신 0.03이었다면” 하는 생각이 들 수 있어요.
하지만 그렇게 하는 순간 저는 텍사스의 명사수가 돼요.
과녁을 먼저 그린 이유는, 결과가 마음에 들지 않을 때도 그 과녁을 그대로 쓰기 위해서예요.

그리고 이 실패는 아무것도 없는 실패가 아니에요.
우리는 꽤 많은 것을 새로 알게 되었어요.

\begin{itemize}
\item 허락을 섞은 동전 랭크는 새 허락을 맞히는 데에서는 뒷면 랭크보다 확실히 나았어요. 봄과 여름 두 번 모두요.
\item 하지만 새 송금자를 맞히는 데에서는 뒷면 랭크만큼 한다고 말할 수 없었어요. 그래서 “모든 면에서 낫다”고는 말할 수 없어요.
\item 동전 랭크를 송금 랭크와 견주면 차이가 더 뚜렷해요. 새 허락에서는 동전 랭크가 0.119 앞섰지만, 새 송금자에서는 0.073 뒤졌어요.
      이 비교도 계획서에 미리 적어 둔 것이에요.
\item 엔도스랭크와 송금 랭크를 여름 기록으로도 채점해 보았어요.
      새 허락은 엔도스랭크(0.333)가 송금 랭크(0.118)보다, 새 송금자는 송금 랭크(0.319)가 엔도스랭크(0.255)보다 잘 맞혔어요. 봄과 같은 모습이에요.
      엔도스랭크는 계획서에 적힌 점수가 아니어서, 등록한 채점이 끝난 뒤에 같은 지갑과 같은 다시 뽑기로 따로 채점했어요.
\item 여름에도 가장 잘 맞힌 것은 받은 허락의 수(새 허락 0.613)와 받은 송금의 수(새 송금자 0.435)였어요.
\end{itemize}

계획서에는 판정에 들어가지 않는 보조 시험도 적어 두었어요.
결과를 숨기지 않으려고, 무엇이든 다 보고하기로 약속했지요.
그리고 뒷면 랭크 대신 송금 랭크를 상대로 세 약속을 다시 채점하는 것,
지도 밖의 지갑을 점수 0이 아니라 화살표 없는 점으로 다루어 다시 채점하는 것도 미리 적어 두었어요.
두 경우 모두 판정은 똑같이 미충족이었어요.

\begin{deeper}[진짜 이름과 숫자]
\begin{itemize}
\item 등록 파일은 \filepath{config/fresh_holdout_2026q3.yaml}, 계획서는 \filepath{docs/fresh_holdout_2026q3_plan.md}이고, 상태를 “등록됨”으로 바꾼 기록 번호(커밋)는 \texttt{4085ab9}예요. 추출 기록에 두 파일의 지문(해시)이 남아 있어요.
\item 6\textasciitilde{}8월 추출은 월별 질의 9개, 과금 약 1.04테라바이트였어요.
\item 논문에서는 약속 1\textasciitilde{}3을 F1\textasciitilde{}F3이라고 불러요. 약속 1은 우월성, 약속 2와 3은 비열등성(허용 폭 0.02) 검정이에요. 모두 짝지은 부트스트랩 400회, 시드 42, 95\% 백분위 구간을 썼어요.
\item 비교 대상인 뒷면 랭크와 앞면 랭크를 등록 문서는 AWP와 EndorseRank라고 불러요(9장).
\item 판정에 들어가지 않는 보조 시험: 보조 시험 1(청산 없이 닫은 비율에서 동전 랭크가 뒷면 랭크보다 나음) 통과, 보조 시험 2(같은 시험에서 동전 랭크가 앞면 랭크보다 나음) 실패, 보조 시험 3(9월 14일 규칙을 새 기간에 다시 적용) 실패.
\item 민감도 분석 두 가지에서도 판정은 같았어요. 송금 랭크를 상대로 하면 약속 1이 $+0.119$ [0.075, 0.169], 약속 2가 $-0.073$ [$-0.092$, $-0.055$], 약속 3이 $+0.037$ [0.011, 0.066]이에요. 다른 하나는 지도 밖 지갑을 고립된 점으로 다룬 것이에요.
\item 등록 채점 뒤에 따로 매긴 엔도스랭크(스펜더 1,964개): 엔도스랭크 $-$ 송금 랭크가 새 허락 $+0.215$ [0.171, 0.270], 송금 랭크 $-$ 엔도스랭크가 새 송금자 $+0.064$ [0.036, 0.088]이에요. 거래자의 청산 없이 닫은 비율은 엔도스랭크 0.081, 송금 랭크 0.045로, 차이 $+0.035$의 범위 [$-0.017$, 0.087]가 0을 넘나들어요.
\end{itemize}
\end{deeper}

\begin{learned}
\begin{itemize}
\item 결과를 보고 나서 찾은 것은 과녁을 나중에 그린 것일 수 있어요.
\item 그래서 계획을 먼저 공개(사전 등록)하고, 아무도 본 적 없는 6\textasciitilde{}8월 기록으로 다시 시험했어요.
\item 세 약속 가운데 약속 1과 약속 3은 통과, 약속 2는 아슬아슬하게 실패해서 판정은 미충족이에요.
\item 그래서 이 논문은 “동전 랭크가 송금 지도만 걷는 것보다 낫다”고 주장하지 않아요. 실패도 그대로 알리는 것이 과학이에요.
\end{itemize}
\end{learned}
